\documentclass[12pt]{article}
\usepackage{mathwizards-carrusel}
\mwcartotal{9}
\begin{document}

% ── 1. Portada ───────────────────────────────────────────────
\mwcarportada{Kit de arranque}{Semana 0: Cositas qué recordar}{@mathwizards.ve}

% ── 2. Álgebra: los 3 errores ────────────────────────────────
\begin{mwcarlamina}{Álgebra: 3 errores que reprueban}
\begin{cajaerror}
  \centering
  $(a+b)^2 = a^2 + b^2$ \quad\mal\\[2pt]
  {\small Correcto: } $(a+b)^2 = a^2 + 2ab + b^2$ \quad\bien
\end{cajaerror}
\begin{cajaerror}
  \centering
  $\sqrt{a^2+b^2} = a + b$ \quad\mal\\[2pt]
  {\small La raíz no se reparte sobre la suma.}
\end{cajaerror}
\begin{cajaerror}
  \centering
  $-(x-3) = -x - 3$ \quad\mal\\[2pt]
  {\small Correcto: } $-(x-3) = -x + 3$ \quad\bien
\end{cajaerror}
\end{mwcarlamina}

% ── 3. Productos notables ────────────────────────────────────
\begin{mwcarlamina}{Productos notables}
\begin{cajateoria}
  \centering
  $(a+b)^2 = a^2 + 2ab + b^2$\\[4pt]
  $(a-b)^2 = a^2 - 2ab + b^2$\\[4pt]
  $(a+b)(a-b) = a^2 - b^2$
\end{cajateoria}
\begin{cajapasos}
  \small El término central $2ab$ es el que más se olvida.
  Siempre: cuadrado del primero, \emph{doble del primero por el
  segundo}, cuadrado del segundo.
\end{cajapasos}
\end{mwcarlamina}

% ── 4. Fracciones algebraicas ────────────────────────────────
\begin{mwcarlamina}{Fracciones: cuándo cancelar}
\begin{cajaok}
  \centering
  $\dfrac{3(x+2)}{3} = x+2$ \quad\bien\\[3pt]
  {\small El 3 multiplica a \emph{todo} el numerador.}
\end{cajaok}
\begin{cajaerror}
  \centering
  $\dfrac{x+2}{x} = 2$ \quad\mal\\[3pt]
  {\small La $x$ está sumando: no es un factor.}
\end{cajaerror}
\begin{cajateoria}
  \centering\nxheav
  Factores sí. Sumandos no.
\end{cajateoria}
\end{mwcarlamina}

% ── 5. Trigonometría mínima ──────────────────────────────────
\begin{mwcarlamina}{Trigonometría mínima}
\begin{cajateoria}
  \centering
  $\sin\theta = \dfrac{\text{opuesto}}{\text{hipotenusa}}$\quad
  $\cos\theta = \dfrac{\text{adyacente}}{\text{hipotenusa}}$\quad
  $\tan\theta = \dfrac{\text{opuesto}}{\text{adyacente}}$
\end{cajateoria}
\begin{cajaformula}
  \centering\small
  \begin{tabular}{c|ccc}
    & $30^\circ$ & $45^\circ$ & $60^\circ$\\\hline
    $\sin$ & $\tfrac12$ & $\tfrac{\sqrt2}{2}$ & $\tfrac{\sqrt3}{2}$\\
    $\cos$ & $\tfrac{\sqrt3}{2}$ & $\tfrac{\sqrt2}{2}$ & $\tfrac12$\\
    $\tan$ & $\tfrac{\sqrt3}{3}$ & $1$ & $\sqrt3$
  \end{tabular}
\end{cajaformula}
\end{mwcarlamina}

% ── 6. Física: unidades ──────────────────────────────────────
\begin{mwcarlamina}{Física: conversión de unidades}
\begin{cajapasos}
  \centering
  $72\ \dfrac{\text{km}}{\text{h}} \times
   \dfrac{1000\ \text{m}}{1\ \text{km}} \times
   \dfrac{1\ \text{h}}{3600\ \text{s}} =
   20\ \dfrac{\text{m}}{\text{s}}$
\end{cajapasos}
\begin{cajateoria}
  \small Las unidades se cancelan con la misma regla:
  \emph{factores sí, sumandos no}. Multiplica por factores que
  valen 1 y deja que las unidades se cancelen.
\end{cajateoria}
\end{mwcarlamina}

% ── 7. Química: el mol ───────────────────────────────────────
\begin{mwcarlamina}{Química: el mol}
\begin{cajateoria}
  \centering
  $1\ \text{mol} = 6{,}022\times 10^{23}$ partículas\\[4pt]
  {\small Es una docena, pero gigante.}
\end{cajateoria}
\begin{cajapasos}
  \centering
  $n = \dfrac{m}{M}$\qquad
  {\small masa $\div$ masa molar $=$ moles}
\end{cajapasos}
\begin{cajaformula}
  \small Todo lo demás del semestre es aplicar factores de conversión.
\end{cajaformula}
\end{mwcarlamina}

% ── 8. Leer una gráfica ──────────────────────────────────────
\begin{mwcarlamina}{Cómo leer una gráfica}
\begin{cajapasos}
  \small
  \textbf{1.} ¿Qué hay en cada eje? (variable y \emph{unidades})\\[3pt]
  \textbf{2.} La pendiente $= \dfrac{\Delta y}{\Delta x}$ te da la
  razón de cambio.\\[3pt]
  \textbf{3.} El intercepto: ¿tiene sentido físico?
\end{cajapasos}
\begin{cajateoria}
  \small Sirve para Cálculo, Física, Química y Bio: una gráfica
  siempre cuenta cómo cambia una cosa respecto a otra.
\end{cajateoria}
\end{mwcarlamina}

% ── 9. Cierre ────────────────────────────────────────────────
\mwcarfin{Guárdalo · Envíaselo a tu sección}{¡Escribe \textbf{ARRANQUE} para enviarte un valioso material de apoyo!}

\end{document}
